\section{Limitations, ethics, and conclusion}
\label{sec:limitations}

The evaluator is synthetic: costs, success rates, disruption processes and penalty weights are not estimated from workers, customers or organizations. No human subjects or personal data are used. Regional AI exposure dispersion is not a fairness guarantee, collaboration rate is not a welfare outcome, and zero disallowed-mode violations is not workplace safety. A poorly chosen real deployment could still automate unequally, overload reviewers, obscure responsibility or displace work. Any field study therefore requires stakeholder participation, privacy governance, domain-specific outcome validity, and human override; NIST and ISO references are framing resources, not certifications~\citep{nist2023airmf,iso42001,iso23894}.

Quantum limitations are equally sharp: four retained ternary variables, exact enumeration, noiseless statevectors, no gate decomposition, no QPU, no noise, no timing and unequal depth-search budgets. Uniform sampling is a controlled distributional baseline, not the strongest possible classical sampler. The cluster-repair method is stronger than a trivial baseline but not a comprehensive solver portfolio. Bootstrap intervals quantify variation across the 32 frozen instances, not generalization to a population of organizations.

Within those limits, Q-OPS demonstrates a disciplined pattern for quantum-assisted organizational research: formulate governance before invocation, make classical controls strong and inspectable, compare distributions on the same retained set and budget, retain adverse results, and keep a classical acceptance boundary. The frozen evidence supports conditional statevector concentration and a lower first-hit draw proxy at $p=2$ relative to uniform, while rejecting claims about monotone depth, universal invocation and complete feasibility. That two-sided result is the paper's central contribution and a reproducible foundation for the open agenda in \cref{sec:agenda}.
